\documentclass[12pt]{article}
\usepackage{geometry}                % See geometry.pdf to learn the layout options. There are lots.
\geometry{letterpaper}                   % ... or a4paper or a5paper or ... 
%\geometry{landscape}                % Activate for for rotated page geometry
%\usepackage[parfill]{parskip}    % Activate to begin paragraphs with an empty line rather than an indent
\usepackage{graphicx}
\usepackage{amssymb}
\usepackage{epstopdf}
\usepackage{hyperref}
\DeclareGraphicsRule{.tif}{png}{.png}{`convert #1 `dirname #1`/`basename #1 .tif`.png}

\title%{An impressionist sketch of gauge-gravity duality}
{Notes on Maldacena duality}
\author{Wm D Linch III}
\date{}                                           % Activate to display a given date or no date


\begin{document}
\maketitle

\tableofcontents

%\abstract
\newpage

\section*{}
{
\begin{flushright}
{\it The goal of this business is to solve QCD.\\
I. Klebanov.}
\end{flushright}
}

%%%%%%%%%%%%%%%%%%%%%%%%%%%%%%%%%%%%%%%%%%%%%%%%%%%
\section{Quantum Field Theory}

%%%%%%%%%%%%%
\subsection{Gauge fields}
Connection $\nabla= \mathrm d+A$ on principle SU($N$)-bundle over 4-dimensional Minkowski space $(\mathbb R^4, \eta)$ with curvature $F=\frac i2[\nabla , \nabla]$.\\
Yang-Mills action $S_\mathrm{YM}\sim \frac 1{g_\mathrm{YM}^2} \int \mathrm d^4 x\, \mathrm {Tr}\,F^2 +\frac \vartheta {4\pi^2} \int F\wedge F $.\\
Symmetries:
\begin{itemize}
\item Global: Scale $x\mapsto \mathrm e^{\Omega_0} x$ in D=4\footnote{Actually, full conformal invariance. The conformal group $\mathrm{Spin}(4,2)\approx \mathrm {SU}(2,2)$ contains the Poicar\'e group $\mathrm{Spin}(3,1)\ltimes \mathbb R^4$ generated by Lorentz transformations and translations $\{ M_{ab}, P_a\}$ extended by conformal and scale transformations $\{K_a, \Delta \}$.}
\item Local: Gauge $\nabla\mapsto g \nabla g^\dagger$ and Weyl $\eta_{ab}\mapsto \mathrm 
e^{2\omega(x)} \eta_{ab}$ in D=4
\item Duality
\end{itemize}
Observables: Gauge invariant combinations $\mathcal O$ of $A$ e.g. $\mathcal O\sim \mathrm{Tr} (F\nabla \nabla ...F...EE...)$ where $E\sim \mathrm P\exp \oint A$.\footnote{The path-ordered exponential $\mathrm P \exp M(t)=1+\int_0^t \mathrm dt^\prime M(t^\prime)+\int_0^t \mathrm dt^\prime\int_0^{t^\prime} \mathrm dt^{\prime\prime} M(t^{\prime}) M(t^{\prime\prime})+\dots$ is the solution to the equation 
\begin{eqnarray}
{\partial \over \partial t} \mathrm P \exp M(t) = M(t) \mathrm P \exp M(t)
\end{eqnarray}
subject to the initial condition $\mathrm P \exp M(0)=1$.
}\\%end footnote
Correlation functions $\langle \mathcal O_1 \dots \mathcal O_n\rangle =  \int \mathcal O_1 \dots \mathcal O_n\,  \mathrm e^{-S_\mathrm{YM}}\,[\mathrm d A]$ e.g. Path integral $Z\sim \int  \mathrm e^{-S_\mathrm{YM}}\,[\mathrm d A]=\langle 1\rangle$\\

{\bf ``Solving QCD''} (pure Yang-Mills in this case) means {\bf find a basis of ``observables'' $\{ \mathcal O\}$ and all correlation functions $\langle \mathcal O\dots \mathcal O\rangle$.}\\

Problems:
\begin{itemize}
\item Gauge symmetry $\Rightarrow$ ghosts
\item UV divergences $\Rightarrow$ Regularization $\Rightarrow$ Renormalization $\Rightarrow$ Running coupling (Scale anomaly) $\Rightarrow$ no perturbation theory at all energies
\item Number of Feynman diagrams at order $n$ is of order $n!(\mathrm{const})^n$ 
\end{itemize}

%%%%%%%%%%%%
\subsubsection{Toy model for quantum effects}
Map $n:\mathbb R^2 \to S^{N-1}$\\
Action for massless, interacting field 
\begin{eqnarray}
S\sim \frac 1{2g_0^2} \int \mathrm d^2 x\sqrt{\gamma}\gamma^{ab} \partial_a n_i(x)\partial_b n^i(x)
~~~:~~~
n_i(x) n^i(x)=1
\end{eqnarray}
Symmetries:
\begin{itemize}
\item Global O($N$) and scale $x\mapsto \mathrm e^{\Omega_0} x$
\item Local Weyl $\gamma_{ab}\mapsto \mathrm e^{2\omega(x)} \gamma_{ab}$
\end{itemize}
Lagrange multiplier field $m(x)$
\begin{eqnarray}
Z\sim \int [\mathrm d n] \mathrm e^{-S} \delta( n^2-1) \sim \int [\mathrm d n]  [\mathrm d m] \mathrm e^{-S-\frac 1{2g_0^2}\int m( n^2-1)}
\end{eqnarray}
makes path integral Gau\ss ian which we can integrate to
\begin{eqnarray}
Z\sim 
	\textrm{``}~\int [\mathrm d m] {g_0^N \over \sqrt {\det^\prime \Delta_m}} \mathrm e^{\frac 1{2g_0^2}\int m} ~ \textrm{''}
\sim \int [\mathrm d m] \mathrm e^{-S^\prime}
\end{eqnarray}
with action $S^\prime = -\frac N2 \int \mathrm d^2 x \left[ -\frac m {g_0^2N}+\log \det^\prime \Delta_m\right]$ and Laplacian $\Delta_m= -\Box +m(x)$.\\
Large $N$ approximation: $N\to \infty$ and $g\to 0$ such that $\lambda_0 = g_0^2 N$ fixed\\
Localization of $Z$ to stationary point of $S^\prime$: $m_*$ s.t. $\left.{\delta S^\prime \over \delta m(x)}\right|_{m=m*}=0$\\
Fourier transform and regulate with regulator $\Lambda$:
\begin{eqnarray}
{\partial \over \partial m} \log {\det}^\prime \Delta_m 
=  {\partial \over \partial m} \log \prod_{p\in \mathbb R^\times} (p^2 + m)
= \int {\mathrm d p\over (2\pi)^2} {1\over  p^2 + m}
=\frac1{4\pi} \log {\Lambda^2 \over m}
\end{eqnarray}
Vacuum expectation value
\begin{eqnarray}
m_*= \Lambda^2 \mathrm e^{-{4\pi\over \lambda_0}}
\end{eqnarray}
implies mass $m^2_n = m_*$ generated (non-perturbatively) for $n(x)$ 
\begin{eqnarray}
S[ n, m=m_*] \sim \frac 1{2g_0^2} \int \mathrm d^2 x \sqrt{\gamma} \left[ \gamma^{ab} \partial_a n_i \partial_b n^i + m^2_n n_i n^i +\mathrm{const.}\right]
\end{eqnarray}
Mass $m_n$ is physical (measured in laboratory) $\Rightarrow$ independent of regularization $\Rightarrow$ coupling not a parameter ${1\over \lambda_0(\Lambda)} = -\frac 1{2\pi} \log {m_n\over \Lambda}$\\
Renormalized coupling: Subtract divergent part ${1\over \lambda(\mu)} = {1\over \lambda_0(\Lambda)} +\frac 1{2\pi} \log {\mu\over \Lambda}$\\

Quantum effects:
\begin{itemize}
\item Scaling and Weyl symmetries anomalous
\item Running coupling: $\mu {\mathrm d\over \mathrm d \mu} \lambda = \beta(\lambda)$ with $\beta$-function $\beta(\lambda)=-\frac1{2\pi} \lambda^2$ %$\frac 1 {\lambda(\mu)} = \frac 1{4\pi} \log {\mu^2\over m^2_n}$
\item Asymptotic freedom: $\lambda \to 0$ as $\mu\to \infty$
\item ... Symmetry restoration
\end{itemize}


%%%%%%%%%%%%%
\subsubsection{Feynman diagrams}
Map $M:\mathbb R^4 \to \mathrm{SL}_N$\\
Action $S\sim \frac 1{g_0^2}  \int \mathrm {Tr} [(\partial M)^2+V(M)]$\\
Partition function $Z\sim \int [\mathrm dM]\,\mathrm e^{-S} $\\
Propagator $\sim g_0^2$\\
Vertices $\sim g_0^{-2}$\\
Loops $\sim N$\\
``Fatgraph'' Feynman diagram with $v$ vertices, $e$ edges, and $f$ faces $\sim (g_0^2)^{e-v}N^f$
\begin{itemize}
\item Euler number $\chi = v-e+f = 2-2g$ on Riemann surface of genus $g$
\item 't Hooft coupling $\lambda = g_0^2 N$
\item Diagrams $(g_0^2)^{e-v}N^f=N^\chi \lambda^{f-\chi}$
\item 't Hooft large $N$ limit: $N\to \infty$ and $g_0\to 0$ with $\lambda$ fixed
\end{itemize}
Effective action localizes on planar diagrams
\begin{eqnarray}
S_\mathrm{eff}=-\log Z\sim \sum_{g=0}^\infty N^{2-2g} f_g(\lambda)\longrightarrow N^2 f_0(\lambda)+O(1)
\end{eqnarray}
When $\lambda \gg 1$, the sum over planar diagrams is dominated by $f\to \infty$ $\Rightarrow$ closed sphere with punctures

%%%%%%%%%%%%%%%%%%%%%%%%%%%%%%%%%%%%%%%%%%%%%%%%%%%
\section{Gravity}

%%%%%%%%%%%%%
\subsection{Strings}
Map $x:(\Sigma, \gamma) \to (\mathbb R^D, \eta)$\\
Polyakov action: 
\begin{eqnarray}
S_x\sim \frac1{2\ell^2} \int \mathrm d^2\sigma \sqrt{\gamma} \gamma^{\mu\nu} \partial_\mu x^a \eta_{ab} \partial_\nu x^b
\end{eqnarray}
Symmetries: 
\begin{itemize}
\item Global target: Poincar\'e: $x \mapsto \Lambda x + a$
\item Local worldsheet: Diff($\Sigma$) and Weyl $\gamma_{\mu \nu}\mapsto \mathrm e^{2\omega(\sigma)} 
\gamma_{\mu \nu}$
\end{itemize}
New term with same symmetries:
\begin{eqnarray}
\chi \sim \frac1{4\pi} \int \mathrm d^2\sigma \sqrt{\gamma}\,R(\gamma)
\end{eqnarray}
Physics determined by symmetries $\Rightarrow S=S_x+\Phi_0 \chi$\\
Partition function:\footnote{Use 1PI effective action instead.}
\begin{eqnarray}
Z\sim \sum_{\Sigma} \int [\mathrm d x] [\mathrm d \gamma] \mathrm e^{-S_x - \Phi_0 \chi}
\sim \sum_{g=0}^\infty (\mathrm e^{-\Phi_0})^{2-2g} f_g(\ell)
\end{eqnarray}
0-loop $\sim \mathrm e^{-2 \Phi_0}$, 1-loop $\sim 1$, ... Closed string loop expansion parameter: $g^2_s=\mathrm e^{2\Phi_0}$\\
{\bf Closed string loop expansion has same form as 't Hooft large-$N$ expansion with}
\begin{eqnarray}
g_s \leftrightarrow \frac 1N \sim g_\mathrm{YM}^2
\end{eqnarray}
Suggests finding string theory such that string tree level $\leftrightarrow$ planar diagrams\footnote{As pointed out to me by Brenno Vallilo, this vision was never fully realized. String theories exhibit open-closed duality which implies $g_c = g_o^2$ (attaching strips to open string worldsheets reduces the Euler number by 1 instead of 2), that is $g_s = g_{YM}^2$ and not $1/N$, as in 't Hooft's large-$N$ expansion. So a purist would say the string expansion is an expansion in $\lambda/N$ and therefore it is not the 't Hooft expansion.}\\
\\
%%%%%%%%%%%%%
%\subsection{Hint's from Liouville}
%Seems we have gauge theory in D=4 $\leftrightarrow$ bosonic strings in D=4\\
String action invariant under diffeomorphisms and local Weyl transformations\\
Use diff to fix metric $\gamma_{ab}\to \mathrm e^{\phi(\sigma)} \delta_{ab}$\\
Requires reparameterization ghosts $(b,c)$\\
Define
\begin{eqnarray}
\mathrm e^{-F[\gamma]}\sim \int [\mathrm d x][\mathrm d b][\mathrm d c]\,\mathrm e^{-S_x[x,\gamma] -S[b,c,\gamma]}
\end{eqnarray}
%\begin{eqnarray}
%\sim \int [\mathrm d x][\mathrm d \phi][\mathrm d b][\mathrm d c]\,\mathrm e^{-S_x[x,g=\delta] -S[b,c,\phi]}
%\end{eqnarray}
Now shift $\gamma \to \mathrm e^\phi(\sigma) \gamma$\\
{\bf Claim:}
\begin{eqnarray}
F[\mathrm e^\phi \gamma]-F[ \gamma] =  {26-D\over 48\pi}\int \mathrm d^2 \sigma \left[ \frac 12 (\partial \phi)^2+\phi R(\gamma) +\mu^2 \mathrm e^\phi\right]
\end{eqnarray}
Liouville action: particle with potential wall $\sim \mathrm e^\phi$\\
Therefore, if we start in D$<26$, we {\em grow a warped extra dimension} \\
\\
Actually, cannot quantize Liouville theory in D$>1$, nevertheless we guess that gauge theory in D=4 corresponds to a string theory in D=5 with a warped extra dimension\\
Most general metric:
\begin{eqnarray}
\label{warped}
\mathrm ds^2 = w^2(z) \left( \mathrm d\mathbf x^2 + \mathrm d z^2\right)
\end{eqnarray}

%%%%%%%%%%%%
\subsection{Backgrounds}
Replace flat background with NLSM
\begin{eqnarray}
S\sim \frac 1{2\ell^2} \int \mathrm d^2\sigma \sqrt{\gamma} \left[ 
\gamma^{\mu \nu} G_{ab}(x) \partial_\mu x^a \partial_\nu x^b 
+\epsilon^{\mu \nu} B_{ab}(x) \partial_\mu x^a \partial_\nu x^b
+ \frac{\ell^2}{2\pi} \Phi(x) R(\gamma)\right]
\end{eqnarray}
Linearized approximation: $G_{ab}(x)= \eta_{ab} + h_{ab}(x)$, $B_{ab}(x)=  b_{ab}(x)$ and $\Phi(x)= \Phi_0+\varphi(x)$\\
Symmetries:
\begin{itemize}
\item Infinitesimal diffeomorphism of target space $\delta h_{ab} =\partial_a \xi_b+ \partial_b \xi_a$ (using equation of motion for $x$)
\item U(1) gerbe relation $\delta b_{ab}=\partial_a \alpha_b - \partial_b \alpha_a$
\end{itemize}
This is the gauge invariance of massless, spin-2 particle {\em suggesting} gravitation on target space\\
Einstein equation $R_{ab}(G)$ for $G$??\\

{\bf figure out how to motivate masslessness, Einstein equation, D=26...}

Bosonic string $\beta$-functions:
\begin{eqnarray}
\beta_{G_{ab}}&=& \ell^2 R_{ab}+ 2 \ell^2\nabla_a\nabla_b \Phi -\frac{\ell^2}4 H_{acd}H_b{}^{cd}+ O(\ell^4)\cr
\beta_{B_{ab}}&=& -\frac{\ell^2}2 \nabla^cH_{abc}+ \ell^2\nabla^c \Phi H_{abc} + O(\ell^4)\cr
\beta_\Phi &=& {D-26\over 6} - {\ell^2\over 2} \nabla^2 \Phi+\ell^2 \nabla\Phi \cdot \nabla \Phi -\frac{\ell^2}{24} H^2+ O(\ell^4)
\end{eqnarray}
The first equation is the Ricci flow, modified by a diffeomorphism. It is the $L^2$ gradient of the 
effective action:
\begin{eqnarray}
S_\mathrm {eff}\sim - \frac1{2{\kappa_0}^2} \int \chi \left[ \Box +R +\frac1{12}H^2\right] \chi+ \dots
\end{eqnarray}
where $\chi^2=\sqrt{G}\mathrm e^{-\Phi}$ is the T-duality invariant dilaton (see below).\\
Local Weyl invariance: $\Phi \mapsto \Phi+D\omega$ and $G_{ab}\mapsto \mathrm e^{2\omega} G_{ab}$\\
String frame: $\omega=0$. Einstein frame: $\Phi=\Phi_0$\\
Gravitational coupling $\kappa=\kappa_0 \mathrm e^{\Phi_0} = \sqrt{8\pi G_N}= \sqrt{8\pi} \ell_\mathrm p^{D-2\over 2}$


%%%%%%%%%%%%%%%%%%%%%
%\subsubsection{An aside on T-duality}
%Consider the compactification of one of the target space directions $x^*$ on a circle of radius $R$. Ignoring the other directions the action is $S\sim \frac 1{\ell^2} \int \mathrm d^2 \sigma G_{**}\partial x^* \bar \partial x^*$. This action is classically equivalent to the action $S\sim \frac 1{\ell^2} \int \mathrm d^2 \sigma \left[G_{**} q^* \bar q^* +i p_* \left(\bar \partial q^* -  \partial \bar q^*\right)\right]$. To see this, we integrate $p_*$ out producing a $\delta$-function in the path integral constraining the ``curl'' of the $q$s. The solution to this constraint is $q^*=\bar \partial x^*$ which, upon substitution, results in the original action. 
%
%On the other hand, the path integral is now quadratic in $q^\prime= q^*+i G^{**}\partial p_*$ 
%$S\sim \frac 1{\ell^2} \int \mathrm d^2 \sigma \left[G_{**} q^\prime \bar q^\prime -G^{**} \partial p_* \bar \partial p_*\right]$
%so it may be integrated to give\footnote{Note that the kinetic term of the dual field has the wrong sign. This is a generic feature of duality and I do not know why it is not addressed in the literature.}
%\begin{eqnarray}
%S\sim \frac 1{\ell^2} \int \mathrm d^2 \sigma \left[-G^{**} \partial p_* \bar \partial p_* - \ell^2 \log G_{**}  \right]
%\end{eqnarray}
%Thus we see that the string theory on the target space with coordinates $x^*$ and metric $G_{**}$ is equivalent to a dual string with coordinates replaced by $x^*\to p_*$ and metric mapped according to the ``Buscher rule''
%\begin{eqnarray}
%G_{**}\mapsto -{1\over G_{**}}
%\end{eqnarray}
%In addition, relaxing the conformal gauge, we see that we have generated a correction to the Fradkin-Tseytlin term $\int \sqrt{\gamma}  R(\gamma)\log G_{**}$. Equivalently, the 1-loop determinant has generated a shift of the dilaton 
%\begin{eqnarray}
%\Phi \mapsto \Phi -  \log G_{**}
%\end{eqnarray}
%(Again, the sign does not agree with that in the literature.) Taken together, we conclude that the combination $\chi$ above is T-duality invariant.

%%%%%%%%%%%%%%%%%%%%%%%%%%%%%%%%%%%%%%%%%%%%%%%%
\section{Putting it all together}
The classical YM theory has the scale invariance $\mathbf x \to \mathrm e^{\Omega_0} \mathbf x$\\
This has to be a symmetry of the string $\Sigma$-model action $\Rightarrow$ symmetry of warped metric (\ref{warped}) $\Rightarrow$ warp factor $w(z) = \frac Lz$ for some constant $L$ setting the length scale:
\begin{eqnarray}
\mathrm d s^2 = L^2 {\mathrm d \mathbf x^2 +\mathrm dz^2\over z^2}
\end{eqnarray}
This is the $AdS_5$ metric in Poincar\'e coordinates. \\
\\
Problems:
\begin{itemize}
\item Yang-Mills theory is not conformal quantum mechanically
\item Easy strings live in D=26 (bosonic) and D=10 (supersymmetric) not D=5
\item $AdS_5$ has $R_{ab}= - (5 L)^{-2} g_{ab} \neq 0$. I.e. not a string background unless $\partial \Phi\neq 0$
\end{itemize}
These problems cancel out:
\begin{itemize}
\item Conformal invariance restored with maximal supersymmetry: $\mathcal N=4$
\item Replace spacetime with $AdS_5 \times X$ where dim($X$)=11 or 5
\item Choose dim($X$)=5 and curvature $R_{ab}|_X = + (5 L)^{-2} g_{ab}$, i.e. $X = S^5$ with radius $L$ {\bf Fix this: Only the scalar curvature vanishes.}
\end{itemize}

{\bf We seek a superstring s.t. the low-energy supergravity on $AdS_5\times S^5$ is dual to $\mathcal N=4$ super-Yang-Mills with gauge group $\mathrm {SU}(N)$ in the large-$N$ limit.}

%%%%%%%%%%%%%%%%
\subsection{Anti-de-Sitter spaces}
Double Wick rotation: $S^n \to H^n \to AdS_n$\\
Homogeneous and isotropic Einstein space:
\begin{eqnarray}
AdS_n\approx Spin(n-1,2)\Big/ Spin(n-1,1)
\end{eqnarray}
Isometries $Spin(4,2)\approx SU(2,2)$ fixing 
\begin{eqnarray}
-X_{-1}^2 - X_0^2 + X_1^2 +\dots +X_{n-1}^2 = L^2
\end{eqnarray}
Poincar\'e coordinates $z^{-1}= L^{-2} (X_{-1} + X_{n-1})$ and $(\mathbf x)_a= z^{-1}L X_a$\\
Induced metric
\begin{eqnarray}
\mathrm d s^2 = {L^2\over z^2}\left( \mathrm d z^2 + \mathrm d\mathbf x^2\right)
\end{eqnarray}



%%%%%%%%%%%%%%%%
\subsection{Maximally supersymmetric Yang-Mills theory}

Super-Poincar\'e algebra $\{P_a, Q_{\alpha A}, \bar Q_{\dot \alpha }^A, M_{ab}, R_A{}^B\}$ with $4\times 4=16$ real supercharges\\
Lorentz:
\begin{eqnarray}
[M , Q]\sim Q ~,~ [M, P]\sim P
\end{eqnarray}
Supersymmetry:
\begin{eqnarray}
\{Q, \bar Q\} \sim P 
\end{eqnarray}
SU(4)$_R$-symmetry:
\begin{eqnarray}
[R, Q]\sim Q 
\end{eqnarray}
Special conformal generator\footnote{The finite action of a special conformal transformation is the combination inversion-translation-inversion:
\begin{eqnarray}
K_a x^b
\sim I P_a I x^b 
\sim I \partial_a {x^b \over x^2}
\sim I {1\over x^2}\left(\delta_a^b -2 {x_a x^b\over x^2}\right) 
\sim x^2\left(\delta_a^b -2 {x_a x^b\over x^2}\right) 
\sim \left(x^2\partial_a -2 x_a x\cdot\partial\right) x^b
\end{eqnarray}
} $K_a\sim x^2 \partial_a - 2x_a x\cdot \partial$ extends:
\begin{eqnarray}
[K, P]\sim M+\Delta
\end{eqnarray}
Dilatation (scale transformation) generator $\Delta \sim x\cdot \partial$:
\begin{eqnarray}
[\Delta , Q]\sim \frac 12 Q ~,~ [\Delta, P]\sim P~,~ [\Delta, K]\sim -K
\end{eqnarray}
Special conformal generators map supercharges to ``$S$-supersymmetry'' charges
\begin{eqnarray}
[K , Q]\sim \bar S ~,~ \{S,\bar S\}\sim K~,~ \{Q, S\} \sim M+\Delta
\end{eqnarray}
Conformal generators double the number of fermionic generators from 16 to 32.\\
Superconformal group PSU(2,2$|$4)$\supset$ SU(2,2)$\times$SU(4) $\approx$ Spin(4,2)$\times$Spin(6):
\begin{eqnarray}
\mathfrak{psu}(2,2|4) &\sim&
\left(\begin{array}{cc}
\begin{array}{cc}
M & P+iK\\
P-iK & \Delta
\end{array}
&Q+i \bar S\\
\bar Q-i S & R
\end{array}\right)
\end{eqnarray}

Field representation (all in adjoint of SU($N$)): 
\begin{eqnarray}
\varphi^{AB} \stackrel{Q_A}{\longrightarrow} \lambda^A\stackrel{Q_A}{\longrightarrow}\mathrm d A
\end{eqnarray}
On-shell supersymmetry $\Rightarrow$ 6 real scalars $\varphi^{AB}=-\varphi^{BA} = \frac12 \epsilon^{ABCD}\bar \varphi_{CD}$\\
In superspace $W^{AB}\sim \varphi^{AB} + \theta^{[A} \lambda^{B]} +\theta^A \sigma^{ab} \theta^B F_{ab}$\\
Action:
\begin{eqnarray}
S\sim \mathrm {Tr}\int \mathrm d^4\theta \,\,\tau \,W^2 +\mathrm{h.c.}
%\left[ \frac1{g^2} F^2 +\frac \vartheta {2\pi} F\tilde F\right]
\end{eqnarray}
Complexified coupling: $\tau = \frac\vartheta{2\pi}+i\frac{2\pi}{g_\mathrm{YM}^2}$\\
Symmetry: $\tau \mapsto \tau+1$ and {\em S-duality} $\tau \mapsto -\frac1\tau$ generate SL(2,$\mathbb Z$)

%%%%%%%%%%%%%%%%%%%%%%%%%%%%%%%%%%%%%%%%%%%%%%%%
\subsection{Type IIB superstrings}
The superstring we seek should have the properties
\begin{itemize}
\item live in D=10 i.e. critical superstring
\item 32 supersymmetries i.e. ``type II''
\item S-duality SL(2,$\mathbb Z$)
\end{itemize}
This string exists and has left- and right-moving spinors of the same chirality\\
IIB Ramond-Ramond field-strengths:
\begin{eqnarray}
\mathbf s \otimes \mathbf s = \bigoplus_{n~\mathrm odd} \underbrace{\mathbf v \wedge \dots \wedge \mathbf v}_{n}&\Rightarrow& F_{[1]},F_{[3]},F_{[5]}^+ 
\end{eqnarray}
Potentials $C_{[0]}$, $C_{[2]}$, $C_{[4]}$ couple to instantons, strings and self-dual 3-branes\\
Together with NS-NS fields $h_{ab}, b_{ab}, \varphi$ $\Rightarrow$ possibility of SL(2,$\mathbb Z$) duality:
\begin{eqnarray}
\tau = \frac\vartheta{2\pi}+i\frac{2\pi}{g_\mathrm{YM}^2} &\leftrightarrow& {{C_{[0]}}\over {2\pi}}+{2\pi}i\mathrm e^{-\Phi_0}
\end{eqnarray}
Consistency check: Are dilaton and RR scalar constant in $AdS_5\times S^5$ background?\\
Action in Einstein frame:
\begin{eqnarray}
S \sim {1\over 2 \kappa^2}\int\mathrm d^{10}x \sqrt{G} \left[ R(G) + \frac12 \sum F^2\right]
\end{eqnarray}
Below we will see that $\int_{S^5} F_{[5]} \sim N$. Then EOM $\Rightarrow$ $R^{10-2}\sim N^{2}$:
\begin{eqnarray}
R\sim (g_\mathrm s N)^{\frac14}\ell_\mathrm s\sim N^{\frac14} \ell_\mathrm p
\end{eqnarray}

%%%%%%%%%%%%%%%%%%%%%%%%%%%%%%%%%%%%%%%%%%%%%%%%
\section{Black branes}
Einstein equation:
\begin{eqnarray}
\mathrm{Ric}-\frac12 g R =8\pi G\, T
\end{eqnarray}
Schwarzschild black hole a.k.a. ``black 0-brane'' has $T_{00}\sim M\delta(r)$:
\begin{eqnarray}
\mathrm d s^2 = -\left( 1-\frac {r_\mathrm S} r\right) \mathrm dt^2 +\left( 1-\frac {r_\mathrm S} r\right)^{-1} \mathrm dr^2+r^2\mathrm d\Omega_2^2
\end{eqnarray}
Schwarzchild radius: $r_\mathrm S=2GM$ with $M$ the gravitational charge $\int T_{00}\, \mathrm {dvol}\sim M$\\
Penrose diagram: \\
Charged Reissner-Nordstr\"om black hole $A_0=Q/r$:
\begin{eqnarray}
\mathrm d s^2 = -\left( 1-\frac {r_+} r\right)\left( 1-\frac {r_-} r\right) \mathrm dt^2 +\left( 1-\frac {r_+} r\right)^{-1}\left( 1-\frac {r_-} r\right)^{-1} \mathrm dr^2+r^2\mathrm d\Omega_2^2
\end{eqnarray}
Roots $r_\pm = \frac12\left(r_\mathrm S\pm \sqrt{r_\mathrm S^2-4r_Q^2}\right)$ of horizon condition $0=g_{00}=1-\frac {r_\mathrm S}r+ \frac {r_Q^2}{r^2}$\\
Charge radius $r_Q = \sqrt{G} Q$ with $Q$ the electromagnetic charge $\int_{S^2} F \sim Q$\\
Cosmic censorship: $r_-\leq r_+ \Rightarrow Q\leq \sqrt{G}M$\\
Extremal black hole (no Hawking radiation) is BPS $Q=M$: ``D0-brane''\\
Near-horizon limit $y=r- r_\mathrm S\to 0$
\begin{eqnarray}
\mathrm d s^2 \to -{y^2\over R^2} \mathrm dt^2 +{R^2\over y^2} \mathrm dy^2+R^2\mathrm d\Omega_2^2
\end{eqnarray}
with $R=2GM=2\sqrt{G}Q$. Poincar\'e coordinates $z=R^2/y$
\begin{eqnarray}
\mathrm d s^2 \to R^2{-\mathrm dt^2 + \mathrm dz^2 \over z^2}+R^2\mathrm d\Omega_2^2
\end{eqnarray}
{\bf Near-horizon geometry of extremal black 0-brane in D=4 is $\mathrm {AdS}_2\times S^2$ with radius $R$ determined by the amount of flux through the 2-sphere.}

%%%%%%%%%%%%%%%%%%%%%%%%%%%%%%%%%%%%%%%%%%%%%%%%%%%%%
\section{Correspondence}
Supergravity approximation $R\gg\ell_s \Rightarrow \lambda \gg 1$ valid for strong gauge theory coupling.\\
Fix radius $R=1 \Rightarrow \ell_\mathrm s^2\sim 1/\sqrt{g_\mathrm s N}$ and $\kappa^{-2}\sim \ell_\mathrm p^{-8}\sim N^{2}$.\\
Gravitational partition function:
\begin{eqnarray}
Z[ \phi \sim \phi_0] = \mathrm e^{-N^2 S_\mathrm {classical} +O(\ell_s^2)}\left( 1+ O(g_s^2)\right)
\end{eqnarray}
Boundary condition $\phi^i(\mathbf x, z) \stackrel{z\to \infty}{\sim}  \phi^i_0(\mathbf x)$ corresponds to an $\mathcal N=4$ super-Yang-Mills operator $\mathcal O_i$.\\
Witten's prescription: $Z[ \phi \sim \phi_0]=\langle \mathrm e^{\int \phi^i_0 \mathcal O_i}\rangle $\\
For large 't Hooft coupling $\lambda \gg 1$ the generating functional for SYM:
\begin{eqnarray}
\langle \mathrm e^{\int \phi^i_0 \mathcal O_i}\rangle
\approx \mathrm e^{-N^2 S_\mathrm {classical}[\phi \sim \phi_0]}
\end{eqnarray}
How to determine $\mathcal O_i$?

%%%%%%%%%%%%%%%%%
\subsection{Example: Scalar field}
Take $\phi$ a scalar field of mass $m$.\\
Action:
\begin{eqnarray}
S\sim N^2 \int \mathrm d^4 x \, \mathrm d z \frac1{z^5} \left[ z^2 (\partial \phi)^2 + m^2 R^2 \phi^2+O(\phi^4)\right] 
\end{eqnarray}
Classical action requires stationary point (EOM):
\begin{eqnarray}
z^3 {\partial\over \partial z} \left( \frac1{z^2} {\partial\phi\over \partial z}\right)-p^2 z^2 \phi - m^2 R^2 \phi= 0
\end{eqnarray}
Near boundary $\phi \sim z^\alpha$:
\begin{eqnarray}
\alpha(\alpha -4) -m^2 R^2=0~\Rightarrow~ \alpha_\pm= 2\pm\sqrt {4+m^2R^2}
\end{eqnarray}
Dominant solution:
\begin{eqnarray}
\phi(x, z) \sim \epsilon^{\alpha_-} \phi_0(x)
\end{eqnarray}
Scaling symmetry $(x, z)\mapsto \mathrm e^{\Omega_0}(x,z) \Rightarrow $ scaling dimension of $\phi_0$ is $\alpha_-$. Then Witten's prescription implies 
\begin{eqnarray}
\Delta_\mathcal O = 2+\sqrt{4+m^2 R^2}
\end{eqnarray}
For dilaton $\Phi$, $m=0$ $\Rightarrow$ $\Delta_\mathcal O = 4$ which suggests $\mathcal O_\Phi$ is the SYM Lagrangian.


%%%%%%%%%%%%%%%%%%%%%%%%%%%%%%%%%%%%%%%%%%%%%%%%%%%%%
%\section{Currents and off-shell representations}
%
%\subsection{Some crap about gauge fields and Noether currents}
%The dynamics of a complex scalar field $\sigma$ is invariant under global phase rotation $\sigma\mapsto \mathrm e^{i \lambda_0} \sigma$. The Noether current $j_a=i \bar \sigma \stackrel{\leftrightarrow}{\partial}_a \sigma$ is conserved on-shell $\Box \sigma =0 $. Note that at this stage we are already using a metric. The coupling of this current to a field is given by $\int A^a j_a$. The conservation of the current implies that the vector is defined only up to the identification $A^a \sim A^a+ \partial^a \alpha$, that is, the current couples to the cohomology class of the (co-)vector field. 
%
%The tensor invariant under the gauge transformation is $F_{ab}=\partial_a A_b - \partial_b A_a$. From this object we can construct a conserved, symmetric tensor $T_{ab}$ provided we use the Maxwell equation $\partial^a F_{ab}=0$ and the Bianchi identity $\epsilon^{abcd} \partial_b F_{cd}=0$. Note that $T_{ab}=F_{ac}F_b{}^c- x \eta_{ab} F^2$ is not conserved unless $x=\frac14$ since otherwise the terms which reconstruct the Bianchi identity do not combine with the same relative weights. Precisely when $x$ takes this value, $T$ is traceless in D=4.
%%kinetic energy $\int F\wedge * F$.\footnote{It is interesting to note that although we needed the full metric to construct the dynamics of $\sigma$, the Weyl part of the metric does not enter the construction of the dynamics of the gauge field. As Leon puts it: the Hodge operator on 2-forms in 4 dimensions is independent of the Weyl component of the metric.}
%
%Proceeding as before, we can couple this tensor to a gauge field as $\int h^{ab} T_{ab}$. By conservation, the gauge field is only defined up to $h^{ab}\sim h^{ab}+\partial^a \xi^b +\partial^b \xi^a$. Furthermore, in D=4 the trace part does not couple $h^{ab}\sim h^{ab}+ \eta^{ab} \omega$. Separating the metric into trace $h$ and traceless $\check h$ parts, the gauge transformations become $\delta \check h_{ab} = (P_1 \xi)_{ab}$ and $\delta h = 2 \partial \cdot \xi + 4 \omega$.\footnote{As usual, $P_1$ denotes the operator mapping vectors to traceless, symmetric rank-2 tensors: $(P_1)_{ab, c} = \partial_{(a} \eta_{b)c} - \frac12 \eta_{ab} \partial_c$.} 
%
%Only the former couples directly to $T$. The natural gauge-invariant tensor for this part is $\check G_{ab}(\check h)= \Box \check h_{ab}-\partial_{(a}\partial^c \check h_{b)c}+\eta_{ab} \partial^c \partial^d \check h_{cd}$. The next step in the sequence is to construct a current $W_{abcd}$ which is conserved modulo the equations of motion for $\check h$. Here something new happens: 
%
%It is conserved $\partial^b \check G_{ab}=0$ but contrary to previous experience, conservation holds off-shell. Additionally it is not traceless $\eta^{ab} \check G_{ab} = 2 \partial^a \partial^b \check h_{ab}$. 
%
%%$\partial_a \check R$ where $\check R=-\partial^a \partial^b \check h_{ab}$.
%
%%%%%%%%%%%%%%%%%%%%
%\subsection{Singletons}
%
%\subsubsection{Dimension 4}
%Isometry group of $AdS_4$: Spin(3,2) $\approx$ Sp(4, $\mathbb R$).\footnote{Remember: this is {\em not} Sp(4) or even Sp(2)!}\\
%Generators of 
%
%\subsubsection{Dimension 5}
%Isometry group of $AdS_5$: Spin(4,2) $\approx$ SU(2,2).
%
%
%%%%%%%%%%%%%%%%%%%%%%%%%%%%%%%%%%%%%%%%%%%%%%%%%%%%%
%\section{Higher-spin}
%
%Everything taken from \cite{Mikhailov}. Oscillator construction of $\mathfrak{su}(2,2)\subset \mathfrak{sp}(8,\mathbb R)$,\footnote{The {\em universal} mathematical convention is that this is (isomorphic to) the 36-dimensional group of $8\times8$ matrices with real entries, preserving the symplectic form. The associated group has $\pi_1 = \mathbb Z$.} acting on symplectic vector ``phase'' space $(\mathbb R^8, \omega)$. Let $J:T\mathbb R^8 \to T\mathbb R^8$ complex structure such that $\omega\in \Omega^{1,1}(\mathbb R^8)$ generated by a quadratic Hamiltonian function $H\in \mathrm{Hom}_\mathbb R (\mathbb R^8\times \mathbb R^8, \mathbb R)$. The collection of all such quadratic Hamiltonians generate the algebra\footnote{We take the complex structure such that it admits holomorphic Lagrangian subspaces.} $\mathfrak{u}(2,2)\subset \mathfrak{sp}(8,\mathbb R)$ with center $\mathfrak u(1)$ generated by $H$.% and commutant $\mathfrak{su}(2,2)$.
%
%The Lagrangian subspace of the complexification of our module is generated by $\{Q^I, \bar Q^I, P_I,\bar P_I\}_{I=1}^4$ which satisfy the ``polarization'' relations 
%\begin{eqnarray}
%[ P_I, \bar Q^J] = \delta^J_I
%\end{eqnarray}
%and conjugate with all other relations trivial. The Hamiltonian may be written in these coordinates as
%\begin{eqnarray}
%H= i\left(\bar P_I Q^I - P_I \bar Q^I  \right)
%\end{eqnarray}
%
%\paragraph{Higher-spin algebra} $\mathfrak {hs}(2,2)$ is the (infinite-dimensional) associative algebra generated by all (not necessarily quadratic) hamiltonians commuting with $H$. Complexify $\mathbb R^8$ and split it into $(1,0)$ and $(0,1)$ parts using the complex structure $J$. Let $\{a^i\}_{i=1}^2$ and $\{\bar a^{\bar \imath}\}_{\bar \imath=1}^2$ generate these spaces. The canonical $U(1)$ symmetry acts with charge $\pm1$ respectively. Then $\mathfrak{hs}(2,2)\subset \mathbb R[a, \bar a]$ is the subalgebra fixed by this action. An element looks like
%\begin{eqnarray}
%f(a,\bar a)=\sum_{I}\alpha_{I\bar I} a^I \bar a^{\bar I}
%\end{eqnarray}
%where $I=(i_1\dots i_I)$ is a multi-index (Weyl-)symmetrized over all permutations. 
%
%Let $I$ denote the ideal generated by ``imaginary'' elements of the form
%\begin{eqnarray}
%f(a,\bar a)+f(-a,\bar a)
%\end{eqnarray}
%Then the subalgebra $\mathfrak{hs}(2,2)/I$ is denoted $\mathfrak {hs}_\mathbb R(2,2)$. Complexifying $\mathfrak {hs}(2,2)\otimes \mathbb C$ introduces the new ideal
%\begin{eqnarray}
%\bar f(a,\bar a)+f(-a,\bar a)
%\end{eqnarray}
%and defines $\mathfrak{hs}_\mathbb C(2,2)$.
%
%
%\paragraph{Fields} $\mathbb F$-valued functions are representations of the $\mathfrak {hs}_\mathbb F(2,2)$ algebra. Let us choose a ``coordinate'' representation such that 
%\begin{eqnarray}
%Q^I f(q, \bar q) = q^I f(q, \bar q) && P_I f(q, \bar q) = {\partial\over \partial q^I} f(q, \bar q) 
%\end{eqnarray}
%{\it et cetera}. The functions invariant under the $U(1)$ action form the, inappropriately named, ``doubleton'' representation of the higher spin algebra $\mathfrak{hs}_\mathbb F(2,2)$ depending on whether $f$ is real of complex as a function.
%
%The doubleton representation admits the inner product 
%\begin{eqnarray}
%(f, g)= \int \overline {f(-q, \bar q)}g(q,\bar q) \mathrm d^2q \mathrm d^2\bar q
%\end{eqnarray}
%Next let $\varphi(x)\in \mathcal S(\mathbb R^4)$ be a {\em spacetime} field (Schwartz function). Consider the ``doubleton transform''
%\begin{eqnarray}
%\hat \varphi(q,\bar q)= \int \mathrm e^{- x_{I\bar J} q^I q^{\bar J}} \varphi(x) \mathrm d^4 x
%\end{eqnarray}
%
%






































\end{document}  

